\chapter{Refactoring, technical debt, legacy systems, and migrations}

\section{Maintenance is the normal state}
Most software cost occurs after the first release: bugs are fixed, dependencies
change, requirements grow, and people learn what the system really does. A design
that is easy to create but hard to change is not finished engineering.

\section{Refactoring}
Refactoring changes internal structure while preserving specified external
behaviour. Fowler's book popularized a catalog of small transformations and the
discipline of keeping the system working while its design improves
\cite{fowler2018refactoring}. The key precondition is a trustworthy safety net:
tests, assertions, domain invariants, or a controlled comparison.

A safe sequence is:

\begin{enumerate}
  \item characterize current behaviour, including surprising behaviour;
  \item make one structural change;
  \item run fast checks and inspect the diff;
  \item preserve or improve the boundary;
  \item repeat, then remove the obsolete path.
\end{enumerate}

Do not call a behaviour change a refactoring. A combined structural and product
change makes failures harder to attribute and rollback more difficult.

\section{Technical debt}
Technical debt is a metaphor for an intentional or accidental future cost. It
can be rational: a temporary adapter may enable a time-critical launch. It is
dangerous when the cost is not recorded, the interest is underestimated, or no
repayment trigger exists. Distinguish debt from mere code you dislike. A messy
but stable module may be cheaper than a fashionable rewrite.

Record debt with:

\begin{itemize}
  \item the shortcut and why it was taken;
  \item the risk or recurring cost;
  \item evidence of impact;
  \item the condition that makes repayment worthwhile;
  \item an owner or review date.
\end{itemize}

\section{Legacy systems}
Legacy means important software whose current behaviour is poorly understood or
whose constraints are hard to change. The first task is not to replace it. It
is to map dependencies, data, users, operational procedures, and hidden
contracts. Add characterization tests at the boundary. Capture current outputs
before changing internals.

A strangler approach can route one capability through a new implementation while
the old system remains. It reduces a big-bang cutover but creates two systems,
data synchronization, and ambiguous ownership during the transition.

\section{Migrations}
For a data migration, define source and target invariants, transformation rules,
handling of invalid rows, reconciliation totals, rollback or compensation, and
the cutover criterion. Run a dry run against representative data. Measure not
only whether rows copied but whether user-visible meaning was preserved.

For a service migration, use shadow traffic only when privacy and side effects
are controlled. Compare outputs with a defined equivalence relation. A new
implementation may be “more correct” and still break clients that depended on
old rounding, ordering, or error codes; decide which compatibility is required.

\section{When to rewrite}
A rewrite is justified by a concrete constraint: unsupported platform, unsafe
security posture, impossible performance target, or cost that exceeds a staged
transition. Estimate the work of reproducing undocumented behaviour, migrating
data, retraining operators, and running both systems. “The code is old” is not a
sufficient business case.

\begin{exercise}
Plan a migration from a legacy order table that stores money as floating point to
integer cents. Include discovery, rounding policy, dual read/write period,
reconciliation, rollback, and the point at which the old column can be removed.
\end{exercise}

\section{Chapter summary}
Maintenance is the main lifecycle, not a cleanup phase. Refactor in small safe
steps, record technical debt as a decision with a trigger, characterize legacy
behaviour, and migrate data with explicit invariants and reconciliation.

